\documentclass[a4paper, landscape, 14pt]{extarticle}
\usepackage[russian]{babel}
\usepackage{amsmath, amssymb}
\usepackage{geometry}

\usepackage{placeins}
\usepackage{graphicx}
\usepackage{listings}
\usepackage{xcolor}

\begin{document}

\title{Сортировки на видеокартах}
\author{Себелев Максим, DeepSeek}
\date{\today}

\maketitle

\newpage

\section*{Введение}
	Поведаю вам про сортировку на gpu. Все, сразу к делу
	\newpage
\section*{Модель памяти OpenCL}

\begin{figure}[!h]
	\centering
	\includegraphics[width=0.8\textwidth]{opencl-memory-model.png}
	
	\label{fig:example}
\end{figure}
\FloatBarrier


\section*{Битонная сортировка}

Битонная сортировка (Bitonic Sort) — это алгоритм сортировки, который крайне хорошо поддается распараллеливанию. Он имеет ассимптотику хуже, чем стандартные алгоритмы сортировки $O(N \log^2 N)$ против $O(N \log N)$, однако при достаточно больших массивах данных (порядка $10^6$ элементов) и грамотном использовании на видеокартах может давать куда лучшие результаты, чем стандартные алгоритмы

\newpage

\section*{Определение:}
	
    \textbf{Битонная последовательность} - это конечная последовательность $\{a_n\}_{n=0}^{N}$, для которой существует k: $\{a_n\}_{n=0}^{k}$ является возрастающей последовательностью, а $\{a_n\}_{n=k + 1}^{N}$ - убывающей (или наоборот). То есть, последовательность является противоположно монотонной слева и справа от своего экстремума.
    \newline
    Для обобщения будем считать монотонные последовательности и последовательности из 1 элемента так же битонными.
   
  	
\section*{Алгоритм реальный}
	\textbf{Небольшая поправка} - алгоритм определен только для массивов с размером равным степени 2. Однако это не проблема, так как засылая данные на видеокарту мы можем добавить элементов до степени 2 элементов. Добавлять будем элементы максимальные по значения для своего типа (например, для {int} - {INT\_MAX}), чтобы после сортировки они остались последними в массиве. В таком случае мы можем считать с видеокарты первые n элементов отсортированного массива, которые будут соотвествовать исходному отсортированному массиву.
	
	Поэтому далее смело рассматриваемым только массивы с длиной равной степени 2.
	\newline
	
	В алгоритме есть 2 понятия: \textbf{step} и \textbf{stage}
	\newline

	Алгоритм состоит из последовательности step-ов, на каждом из которых выполняется определнная последовательность stage-й.
	\newline
	
	\textbf{step:}
	На i-ом step-е последовательно выполняем stage-и от i-го до 0-го.
	\newline

	\textbf{stage:}
	На i-ом stage делим массив на последовательные боксы по $2^{i+1}$ элементов (именно из-за этого шага длина массива должна быть равна степени 2). Элементы внутри этого бокса пронумеруем от 0 до $2^{i+1}$. Элементы с разницей индексов равной $2^{i}$ выстраиваем их в отсортированном порядке. Именно этот шаг и поддается распараллеливанию, потому что все элементы делятся по парам и этим шаги сортировки внутри бокса выполняются для таких пар независимо и так же для всех боксов эти вычисления проводятся независимо.

\begin{figure}[!h]
	\centering
	\includegraphics[width=1\textwidth]{bitonic-sort.png}

	\label{fig:example}
\end{figure}
\FloatBarrier


\section*{А зачем было определение? И где битонность то?}

	Алгоритм, описанный выше на самом деле крайне незаметно создает битонные последовательности и так же незаметно сливает и сплитит х:
	Действительно, после i-го step-а массив делится на i+1 битонных последовательностей. А на следующем step-е соседние последовательности (0 и 1, 2 и 3, 4 и 5, ...) сливаются в одну большую.
	Собственно здесь замаскировано простым алгоритмом, что серьезные ученые называют bitonic-split и bitonic-merge, отсюда и название. Однако заумные слова не помогают написать эффектичный код на OpenCL, поэтому не будем заострять внимание на этом.
	
	\begin{figure}[!h]
		\centering
		\includegraphics[width=1\textwidth]{bitonic-sorting-net.png}
		
		\label{fig:example}
	\end{figure}
	\FloatBarrier
\newpage

\section*{Реализация}
\section*{Первые шаги}
	Давайте заметим, что размер локальной памяти для CU обычно довольно маленький. Значит не все боксы мы сможем туда помещать. Поэтому в локальной памяти мы сможет работать только на небольших stage-ах. На моем устройстве размера локальной памяти CU хватит лишь на 512 4-байтных int.
	Теперь давайте заметим, что для первых step-ов все stage будут влезать в локальную память, поэтому вынесем в отдельное ядро функцию, отвечающие за первые step-ы:

	
	\begin{figure}[h]
		\centering
		\includegraphics[width=1.1\textwidth]{first-steps.png}
		
		\label{fig:example}
	\end{figure}
\FloatBarrier

\section*{Шире шаг!}
	
	Теперь давайте заметим, что даже на больших step-ах у нас есть stage-и в которых все боксы помещающтся в локальную память. Поэтому разделим большие step-ы на 2 этапа:
	\begin{enumerate}
		\item выполняем большие stage в глобальной памяти
		\item выполняем маленькие stage в локальной памяти
	\end{enumerate}

	Итерацию же по step и stage мы вынуждены вынести на cpu, так как барьеры могут быть только по локальной памяти (особенность gpu), а нам нужна синхронизнация после каждого step на всем массиве.
	
	\newpage
	\section*{Шаг широкий, глобально}
	
	\begin{figure}[!h]
		
		\noindent\includegraphics[width=1.3\textwidth]{big-steps-big-stages.png}
		
		\label{fig:example}
	\end{figure}
	\FloatBarrier
	
	\newpage
	\section*{Шаг широкий, локально}
	\begin{figure}[!h]
		
		\noindent\includegraphics[width=0.9\textwidth]{big-steps-small-stages.png}
		
		\label{fig:example}
	\end{figure}
	\FloatBarrier

	\newpage

	\section*{Дошагали до генштаба}
	\begin{figure}[!h]
		\noindent\includegraphics[width=1.1\textwidth]{c++-driver.png}
		
		\label{fig:example}
	\end{figure}
	\FloatBarrier

	\newpage
	\section*{И чего шагали, спрашивается (aka бенчмарк)}

	\textbf{Что за видеокарта, что за процессор?}
	\newline
	Все бенчмарки проводились на моем старичке \textbf{Lenovo ThinkPad X1 Carbon Gen 9}, 21 года рождения,
	\newline
	с процессором
	\textbf{11th Gen Intel® Core™ i7-1185G7 × 8},
	\newline
	и встроенной видеокартой
	\textbf{Intel Iris XE Graphics}
	
	*танки еле тянет на минимальных настройках, поэтому достигнутые результаты, о которых далее, должны быть подвергнуты удвоенному восхищению*
	\newline



	Посмотрим результаты бенчмарка (все время в милисекундах):
	
	\newpage
	
	
	\begin{figure}[!h]		
		
		\noindent\includegraphics[width=1.2\textwidth]{benchmark-stdsort-win.png}
		
		\label{fig:example}
	\end{figure}
	\FloatBarrier
	
		На маленьких массивах, как это и ожидалось, накладные расходы от копирования данных на видеокарту и прочей работы с видеокартой перевешивают эффективность паралленьного выполнения сортировки, и потому std::sort с большим отрывом побеждает. Однако...

	\newpage
	\begin{figure}[!h]		
		
		\noindent\includegraphics[width=1.3\textwidth]{benchmark-bitonicsort-win.png}
		
		\label{fig:example}
	\end{figure}
	\FloatBarrier
		... когда мы приближаемся к полумиллиону элементов, преимущество std::sort начинает невелироваться, а начиная с 600\_000 элементов bitonicsort начинает обгонять std::sort, доводя свое преимущество до 5-кратного на 100\_000\_000  элементов:
		
		\newpage
		
		\centering\section*{ну все, покедова}
\vspace{0.5cm}


\end{document}


	
